\subsection{Enumeración por Bitmask (bits perdidos de un mensaje)}
\label{alg:5-2}

\graybold{Preguntas guía:}

\begin{itemize}
\item ¿Hay un pequeño número (≤ \aprox{}20) de decisiones binarias desconocidas (bits, elecciones sí/no) y hace falta probar combinaciones?
\item ¿El espacio total de combinaciones (\pow{2}{k}) es lo bastante chico para probarlas todas dentro del límite de tiempo?
\item ¿Cada combinación se puede armar y evaluar rápido a partir de su representación binaria?
\end{itemize}

\graybold{Utilidad:} Representar y recorrer exhaustivamente todas las configuraciones posibles de un conjunto pequeño de decisiones binarias. Cada máscara de bits codifica una asignación distinta de valores, permitiendo generar candidatos de forma sistemática y verificar cuál satisface la condición del problema.

\graybold{Complejidad:} \bigO{\pow{2}{k} \cdot C}, con $k$ = número de incógnitas binarias y $C$ = costo de evaluar una combinación.

\graybold{Fundamento teórico:} Cualquier asignación de k bits desconocidos corresponde biunívocamente a un entero entre 0 y $\pow{2}{k}-1$. Recorriendo ese rango y extrayendo el bit j con (mascara \textgreater{}\textgreater{} j) \& 1, se reconstruye cada combinación posible sin necesidad de recursión ni backtracking explícito. Es la enumeración más directa de un espacio de búsqueda combinatorio pequeño.

\graybold{Ejercicio aplicado}

Un mensaje M' y un divisor N' se recibieron con hasta 16 bits en total marcados como perdidos ('*'). Reconstruir un mensaje M compatible (M completando los '*' de M') tal que exista un N (completando los '*' de N') donde N divide a M exactamente.

\graybold{I/O:} |M'| ≤ 500 pero el total de '*' (en M' y N' combinados) es ≤ 16 → lectura total con readline es más que suficiente; el costo real está en las \pow{2}{16} combinaciones, no en la entrada.

\graybold{Código solución}

\begin{lstlisting}[language=Python]
import sys

input = sys.stdin.readline


def generate(strings):
    pos = []            # posiciones donde existen bits desconocidos
    for i, c in enumerate(strings):
        if c == '*':
            pos.append(i)
    k = len(pos)

    ans = []
    for mask in range(1 << k):          # cada mascara representa una asignacion distinta de los k bits perdidos
        cur = strings[:]
        for j in range(k):
            if (mask >> j) & 1:
                cur[pos[j]] = '1'
            else:
                cur[pos[j]] = '0'
        ans.append(''.join(cur))
    return ans

M = list(input().strip())
N = list(input().strip())

messages = generate(M)
divisors = generate(N)

for m in messages:              # probar todas las combinaciones posibles
    value_m = int(m, 2)
    for n in divisors:
        value_n = int(n, 2)

        if value_n == 0:
            continue            # la division entre cero no existe
        if value_m % value_n == 0:        # la primera combinacion valida es respuesta
            print(m)
            sys.exit()
\end{lstlisting}
